
%(BEGIN_QUESTION)
% Copyright 2006, Tony R. Kuphaldt, released under the Creative Commons Attribution License (v 1.0)
% This means you may do almost anything with this work of mine, so long as you give me proper credit

En interessant måte å skrive spektroskopisk elektronskallstruktur for grunnstoffer på er vist her:

% No blank lines allowed between lines of an \halign structure!
% I use comments (%) instead, so that TeX doesn't choke.

$$\vbox{\offinterlineskip
\halign{\strut
\vrule \quad\hfil # \ \hfil & 
\vrule \quad\hfil # \ \hfil \vrule \cr
\noalign{\hrule}
%
% First row
Element & Elektronkonfigurasjon \cr
%
\noalign{\hrule}
%
% Another row
Hydrogen & 1s$^{1}$ \cr
%
\noalign{\hrule}
%
% Another row
Helium & 1s$^{2}$ \cr
%
\noalign{\hrule}
%
% Another row
Lithium & [He]2s$^{1}$ \cr
%
\noalign{\hrule}
%
% Another row
Beryllium & [He]2s$^{2}$ \cr
%
\noalign{\hrule}
%
% Another row
Boron & [He]2s$^{2}$2p$^{1}$ \cr
%
\noalign{\hrule}
%
% Another row
Carbon & [He]2s$^{2}$2p$^{2}$ \cr
%
\noalign{\hrule}
%
% Another row
Nitrogen & [He]2s$^{2}$2p$^{3}$ \cr
%
\noalign{\hrule}
%
% Another row
Oxygen & [He]2s$^{2}$2p$^{4}$ \cr
%
\noalign{\hrule}
%
% Another row
Fluorine & [He]2s$^{2}$2p$^{5}$ \cr
%
\noalign{\hrule}
%
% Another row
Neon & [He]2s$^{2}$2p$^{6}$ \cr
%
\noalign{\hrule}
%
% Another row
Sodium & [Ne]3s$^{1}$ \cr
%
\noalign{\hrule}
%
% Another row
Magnesium & [Ne]3s$^{2}$ \cr
%
\noalign{\hrule}
%
% Another row
Aluminum & [Ne]3s$^{2}$3p$^{1}$ \cr
%
\noalign{\hrule}
%
% Another row
Silicon & [Ne]3s$^{2}$3p$^{2}$ \cr
%
\noalign{\hrule}
%
% Another row
Phosphorus & [Ne]3s$^{2}$3p$^{3}$ \cr
%
\noalign{\hrule}
%
% Another row
Sulfur & [Ne]3s$^{2}$3p$^{4}$ \cr
%
\noalign{\hrule}
%
% Another row
Chlorine & [Ne]3s$^{2}$3p$^{5}$ \cr
%
\noalign{\hrule}
%
% Another row
Argon & [Ne]3s$^{2}$3p$^{6}$ \cr
%
\noalign{\hrule}
%
% Another row
Potassium & [Ar]4s$^{1}$ \cr
%
\noalign{\hrule}
} % End of \halign 
}$$ % End of \vbox

Forklar hvordan denne notasjonen komprimerer et ellers langt og tungvint format (der {\it alle} skall i hvert atom vises), uten å utelate avgjørende informasjon.

\vskip 20pt \vbox{\hrule \hbox{\strut \vrule{} {\bf Forslag til sokratisk diskusjon} \vrule} \hrule}

\begin{itemize}
\item{} Grunnstoffnavnet i hakeparentes omtales noen ganger som en ``electron core'' i spektroskopisk notasjon. Forklar hvorfor dette er en fornuftig betegnelse.
\item{} Se på tabellen, og identifiser ``periodene'' som utgjør hver rad i periodesystemet.
\item{} Identifiser grunnstoffene i tabellen som er lettest å ionisere positivt (fjerne ett elektron).
\item{} Identifiser grunnstoffene i tabellen som er lettest å ionisere negativt (ta opp ett elektron).
\end{itemize}

\underbar{file i00562}
%(END_QUESTION)





%(BEGIN_ANSWER)


%(END_ANSWER)





%(BEGIN_NOTES)

Denne (forkortede) spektroskopiske notasjonen er enklere enn fullversjonen fordi repeterende deler ``komprimeres'' til neste lavere grunnstoffsymbol med fullt fylte elektronskall. For eksempel gjentas elektronkonfigurasjonen til neon (1s$^{2}$2s$^{2}$2p$^{6}$) i alle grunnstoffer med atomnummer over 10, så i stedet for å skrive hele elektronrekkefølgen for alle høyere grunnstoffer representeres 1.- og 2.-skallet av symbolet for neon [Ne].

\vskip 10pt

En kjemiker kan si at elektronkonfigurasjonen til natrium for eksempel er ett ``s''-subskallelektron i 3. skall (3s$^{1}$), over en neon-{\it core}.

Den {\it virkelig} korte notasjonen som finnes i de fleste periodesystem viser bare ytterste skall, med en implisitt ``core'' av lavere fylte skall som ikke vises.

%INDEX% Fysikk, atom: elektronskall (spektroskopisk notasjon)
%INDEX% Kjemi, grunnprinsipper: elektronskall (spektroskopisk notasjon)

%(END_NOTES)

